\documentclass[aspectratio=169]{beamer}

\usepackage[utf8]{inputenc}
\usepackage[T1]{fontenc}
\usepackage{lmodern}
\usepackage{booktabs}
\usepackage{graphicx}
\usepackage{xcolor}
\usepackage{hyperref}

\usetheme{Madrid}
\usecolortheme{seagull}

\title{Vers une approche adaptative combinant IA générative et gamification}
\subtitle{Apprentissage immersif et personnalisé dans les environnements e-learning}
\author{Zouhair El Hadiq}
\institute{Professeur agrégé d'informatique -- CPGE Bab Sahra (Guelmim)\\Fondateur de la plateforme \texttt{cpgeacademy.org}}
\date{Audition doctorale avec Pr. Zahi Jarir -- 2025}

\begin{document}

\begin{frame}
  \titlepage
\end{frame}

\begin{frame}{Plan de l'exposé}
  \tableofcontents
\end{frame}

\section{Contexte et ambition}

\begin{frame}{Pourquoi ce projet ?}
  \begin{itemize}
    \item Taux d'abandon élevé (70--80 \%) dans les formations e-learning malgré l'abondance de contenus.
    \item Hétérogénéité forte des profils CPGE et professionnels : besoin de personnalisation fine.
    \item Potentiel des LLM pour générer contenus et feedbacks contextualisés, encore peu exploité pédagogiquement.
    \item Volonté du laboratoire de consolider l'axe \textbf{IA éducative} avec un terrain réel (CPGE, formation continue).
  \end{itemize}
\end{frame}

\begin{frame}{Alignement profil-sujet}
  \begin{columns}
    \column{0.55\textwidth}
    \begin{itemize}
      \item Professeur agrégé d'informatique à la CPGE Bab Sahra (Guelmim), expérience terrain MP/PSI.
      \item Ingénieur en modélisation et calcul scientifique (EMI Rabat).
      \item Créateur et pilote de \texttt{cpgeacademy.org} (ressources numériques, analytics).
      \item Encadrement TIPE, productions de contenus adaptatifs, animation d'équipes pédagogiques.
    \end{itemize}
    \column{0.45\textwidth}
    \begin{block}{Atout principal}
      Capacité à \textbf{relier recherche et pratique} : expérimentation rapide, retours d'apprenants, diffusion.
    \end{block}
  \end{columns}
\end{frame}

\section{Problématique et état de l'art}

\begin{frame}{Question de recherche}
  \begin{block}{Formulation}
    Comment concevoir un système e-learning qui s'adapte dynamiquement au profil \textbf{cognitif} et \textbf{motivationnel} de chaque apprenant en combinant \textbf{IA générative} pour la création de contenus personnalisés et \textbf{gamification} pour soutenir l'engagement ?
  \end{block}
  \begin{itemize}
    \item Cible : environnements e-learning intensifs (CPGE, formation continue).
    \item Objectif : immersion, persévérance, amélioration mesurable des performances.
  \end{itemize}
\end{frame}

\begin{frame}{État de l'art synthétique}
  \begin{columns}
    \column{0.33\textwidth}
    \textbf{IA générative}
    \begin{itemize}
      \item ChatGPT, Gemini, LLaMA pour contenus/feedbacks.
      \item Tuteurs conversationnels (Luckin, 2024).
      \item Limites : hallucinations, alignement pédagogique.
    \end{itemize}
    \column{0.33\textwidth}
    \textbf{Gamification}
    \begin{itemize}
      \item Self-Determination Theory (Deci \& Ryan).
      \item Mécaniques : quêtes, badges, progression.
      \item Risque : sur-justification extrinsèque.
    \end{itemize}
    \column{0.33\textwidth}
    \textbf{Adaptativité}
    \begin{itemize}
      \item Modèles d'apprenants (IRT, BKT).
      \item Plateformes ALEKS, Knewton.
      \item Faible intégration avec LLM temps réel.
    \end{itemize}
  \end{columns}
\end{frame}

\section{Approche proposée}

\begin{frame}{Architecture prévisionnelle}
  \begin{itemize}
    \item \textbf{Profil apprenant} multidimensionnel : connaissances, style, motivation, progression.
    \item \textbf{Moteur IA générative} encadré :
          \begin{itemize}
             \item prompts structurés, corpus validé CPGE.
             \item approche RAG avec base de connaissances disciplinaires.
             \item filtres pédagogiques et relecture automatisée.
          \end{itemize}
    \item \textbf{Moteur de gamification adaptatif} :
          \begin{itemize}
            \item quêtes et défis alignés sur la progression réelle.
            \item feedback instantané, indicateurs d'engagement (flow, temps, assiduité).
          \end{itemize}
    \item \textbf{Couches d'analytics} : tableau de bord enseignants, diagnostics apprenants, suivi longitudinal.
  \end{itemize}
\end{frame}

\begin{frame}{Feuille de route (36 mois)}
  \begin{center}
    \begin{tabular}{@{}p{0.27\textwidth}p{0.3\textwidth}p{0.32\textwidth}@{}}
      \toprule
      Phase & Objectifs clés & Livrables \\
      \midrule
      Exploration (0--9 mois) & Revue, modèle conceptuel, protocole éthique & Cartographie de l'état de l'art, schéma de données, grille d'évaluation \\
      Prototypage (9--24 mois) & Développement plateforme, intégration LLM, moteur gamifié & Prototype fonctionnel, prompts pédagogiques, scénarios adaptatifs \\
      Validation (24--36 mois) & Expérimentations CPGE et formation pro & Jeux de données anonymisés, rapports d'études, articles soumis \\
      \bottomrule
    \end{tabular}
  \end{center}
\end{frame}

\begin{frame}{Méthodologie d'évaluation}
  \begin{itemize}
    \item \textbf{Design expérimental} : groupes contrôle/expérimental, études longitudinales.
    \item \textbf{Mesures d'apprentissage} : pré/post-tests, progression sur exercices, indicateurs de transfert.
    \item \textbf{Mesures d'engagement} : temps actif, taux de complétion, score de flow, auto-évaluations.
    \item \textbf{Analyse qualitative} : entretiens semi-directifs, feedback enseignants, traces d'interaction.
    \item \textbf{Aspects éthiques} : consentement, anonymisation, politique d'usage IA responsable.
  \end{itemize}
\end{frame}

\section{Valeur ajoutée et défis}

\begin{frame}{Contributions attendues}
  \begin{columns}
    \column{0.5\textwidth}
    \textbf{Scientifiques}
    \begin{itemize}
      \item Cadre théorique unifiant IA générative et gamification adaptative.
      \item Modèle de profil apprenant intégrant dimensions motivationnelles.
      \item Guides méthodologiques pour l'implémentation d'IA générative en pédagogie.
    \end{itemize}
    \column{0.5\textwidth}
    \textbf{Opérationnelles}
    \begin{itemize}
      \item Extension immédiate sur \texttt{cpgeacademy.org}.
      \item Démos pour formations pro et partenaires industriels.
      \item Outils d'analyse pour enseignants CPGE.
    \end{itemize}
  \end{columns}
\end{frame}

\begin{frame}{Pourquoi moi ?}
  \begin{itemize}
    \item Expérience unique avec des cohortes à forte intensité cognitive (CPGE) et retours d'apprenants.
    \item Capacité prouvée à produire et déployer des ressources numériques (plateforme, analytics).
    \item Maîtrise technique (IA, modélisation) et pédagogique (progression CPGE, TIPE).
    \item Accès à un terrain d'expérimentation immédiatement exploitable et volonté de diffusion.
    \item Motivation forte à s'engager trois ans pour faire rayonner le laboratoire.
  \end{itemize}
\end{frame}

\begin{frame}{Défis anticipés et réponses}
  \begin{itemize}
    \item \textbf{Technique} : coût et latence des LLM \textrightarrow{} optimisation, cache, modèles hybrides.
    \item \textbf{Pédagogique} : garantir la rigueur du contenu \textrightarrow{} validation par enseignants, référentiels CPGE.
    \item \textbf{Éthique} : biais, dépendance \textrightarrow{} chartes d'usage, suivi humain, transparence.
    \item \textbf{Scientifique} : mesure de l'effet à long terme \textrightarrow{} protocoles longitudinal et mix-methods.
  \end{itemize}
\end{frame}

\section{Conclusion}

\begin{frame}{Feuille de route immédiate}
  \begin{itemize}
    \item Affiner la revue bibliographique et sélectionner 3 cas d'étude alignés avec le laboratoire.
    \item Co-construire le cahier des charges du prototype avec l'équipe encadrante.
    \item Préparer une expérimentation pilote sur un module cpgeacademy (été 2025).
    \item Définir les modalités de suivi (comité, livrables, jalons semestriels).
  \end{itemize}
\end{frame}

\begin{frame}{Conclusion}
  \begin{itemize}
    \item Sujet proposé par le Pr. Jarir qui s'appuie sur un besoin identifié sur le terrain CPGE.
    \item Candidat disposant du triptyque clé \textbf{(expertise pédagogique, maîtrise IA, terrain d'expérimentation)}.
    \item Ambition : construire un système immersif et personnalisable qui améliore durablement la réussite en e-learning.
  \end{itemize}
  \vfill
  \begin{block}{Ouverture}
    Je propose de lancer un atelier commun laboratoire--CPGE pour co-désigner le prototype et définir ensemble les premiers cas d'usage.
  \end{block}
\end{frame}

\end{document}
